\documentclass[11pt]{article}

\usepackage[margin=1in]{geometry}
\usepackage[T1]{fontenc}
\usepackage{lmodern}
\usepackage{amsmath,amssymb,amsthm,mathtools}
\usepackage{microtype}
\usepackage[colorlinks=true,linkcolor=blue,urlcolor=blue]{hyperref}

\newtheorem{theorem}{Theorem}
\newtheorem{corollary}[theorem]{Corollary}
\newtheorem{remark}[theorem]{Remark}
\newcommand{\R}{\mathbb{R}}
\newcommand{\C}{\mathbb{C}}
\newcommand{\A}{\mathcal{A}}
\newcommand{\dd}{\,\mathrm{d}}

\title{Whole-Space Nonattainment and Strict Spreading for a Positive Curvature Functional}
\author{Richard J. Reyes}
\date{October 26, 2026}

\begin{document}
\maketitle

\begin{abstract}
For a nonnegative energy comprising a Dirichlet term and a bounded,
nonlinear weight multiplying the squared Laplacian, we show that the
fixed-positive-mass infimum on $H^2(\R^d;\C)$ is zero and is not attained
in any finite spatial dimension $d\ge 1$. For $1\le d\le4$, the energy
is strictly increasing along the mass-preserving dilation parameter,
so no positive-mass local constrained minimizer exists in the strong
$H^2$ topology. The conclusions are variational and do not exclude
stationary saddle points or localization supplied by other terms or
boundary conditions.
\end{abstract}

\section{Functional and results}
Fix $d\in\mathbb{N}$ and constants $c_1,c_2,\varepsilon>0$ and
$\alpha\ge0$. For $s\ge0$ define
\begin{equation}\label{eq:weight}
 F(s)=\frac{s}{\bigl(s+\varepsilon^2e^{-2\alpha s}\bigr)^2}.
\end{equation}
For $\psi\in H^2(\R^d;\C)$, set
\begin{equation}\label{eq:energy}
 E(\psi)
 =c_1\int_{\R^d}|\nabla\psi|^2\dd x
 +c_2\int_{\R^d}F(|\psi|^2)|\Delta\psi|^2\dd x.
\end{equation}
For every mass $m>0$ define
\begin{equation}\label{eq:mass}
 \A_m=\bigl\{\psi\in H^2(\R^d;\C):\|\psi\|_{L^2}^2=m\bigr\},
 \qquad I_d(m)=\inf_{\psi\in\A_m} E(\psi).
\end{equation}

\begin{theorem}[Positive-mass whole-space nonattainment]\label{thm:global}
For every finite dimension $d\ge1$ and every $m>0$, the energy
\eqref{eq:energy} is finite and nonnegative on $H^2(\R^d;\C)$,
strictly positive for each nonzero $\psi\in H^2(\R^d;\C)$, and
\begin{equation}\label{eq:main}
 \boxed{I_d(m)=0,\qquad
 \operatorname*{arg\,min}_{\psi\in\A_m}E(\psi)=\varnothing.}
\end{equation}
\end{theorem}

\begin{theorem}[Strict spreading for dimensions one through four]\label{thm:monotone}
Suppose $1\le d\le4$. For each nonzero
$\psi\in H^2(\R^d;\C)$ and $\lambda>0$, let
\begin{equation}\label{eq:dilation}
 \psi_\lambda(x)=\lambda^{d/2}\psi(\lambda x).
\end{equation}
The function $\lambda\mapsto E(\psi_\lambda)$ is continuously
differentiable and strictly increasing on $(0,\infty)$. In particular,
\begin{equation}\label{eq:strict}
 0<\lambda<1\quad\Longrightarrow\quad E(\psi_\lambda)<E(\psi).
\end{equation}
Consequently, $E|_{\A_m}$ has no local minimizer in the $H^2$
topology for any $m>0$.
\end{theorem}

\section{Proof of all-dimensional nonattainment}
\begin{proof}[Proof of Theorem~\ref{thm:global}]
Write
\begin{equation*}
 A(s)=\varepsilon^2e^{-2\alpha s},\qquad D(s)=s+A(s).
\end{equation*}
Since $A(s)>0$, the function $F(s)=s/D(s)^2$ is continuous
and nonnegative on $[0,\infty)$, and positive when $s>0$.
Moreover,
\begin{equation}\label{eq:asymptotics}
 F(s)=\frac{s}{\varepsilon^4}+O(s^2)\quad(s\downarrow0),
 \qquad F(s)=\frac{1+o(1)}{s}\quad(s\to\infty).
\end{equation}
It follows that $F$ is bounded; set
$M_F=\sup_{s\ge0}F(s)<\infty$.
For every $\psi\in H^2$,
\begin{equation}\label{eq:bound}
 0\le \int_{\R^d}F(|\psi|^2)|\Delta\psi|^2\dd x
 \le M_F\|\Delta\psi\|_{L^2}^2<\infty.
\end{equation}
Thus the energy is finite and nonnegative in all dimensions.

Furthermore, if $E(\psi)=0$, then $c_1>0$ implies
$\nabla\psi=0$ almost everywhere. Since $\R^d$ is connected,
$\psi$ is almost everywhere constant. The only constant in
$L^2(\R^d)$ is zero; hence
\begin{equation}\label{eq:positive}
 E(\psi)=0\quad\Longleftrightarrow\quad\psi=0.
\end{equation}

Choose any $\psi\in\A_m$. For $\lambda>0$ use the dilation
\eqref{eq:dilation}; a change of variables gives
\begin{equation}\label{eq:mass_preserved}
 \|\psi_\lambda\|_{L^2}^2=\|\psi\|_{L^2}^2=m,
 \qquad \|\nabla\psi_\lambda\|_{L^2}^2
 =\lambda^2\|\nabla\psi\|_{L^2}^2.
\end{equation}
Also $\Delta\psi_\lambda(x)=\lambda^{d/2+2}\Delta\psi(\lambda x)$.
Writing $\rho=|\psi|^2$ and $B=|\Delta\psi|^2$, we obtain the
exact scaling formula
\begin{equation}\label{eq:scale}
 E(\psi_\lambda)
 =c_1\lambda^2\|\nabla\psi\|_{L^2}^2
 +c_2\lambda^4\int_{\R^d}F(\lambda^d\rho(y))B(y)\dd y.
\end{equation}
Using the uniform bound on $F$,
\begin{equation}\label{eq:upper}
 0\le E(\psi_\lambda)
 \le c_1\lambda^2\|\nabla\psi\|_{L^2}^2
 +c_2 M_F\lambda^4\|\Delta\psi\|_{L^2}^2
 \xrightarrow[\lambda\downarrow0]{}0.
\end{equation}
Since $\psi_\lambda\in\A_m$, nonnegativity and
\eqref{eq:upper} show $I_d(m)=0$.
If some $\psi_*\in\A_m$ attained this infimum, then
$E(\psi_*)=0$ and \eqref{eq:positive} would force
$\psi_*=0$, contradicting $\|\psi_*\|_{L^2}^2=m>0$.
Hence the infimum is not attained.
\end{proof}

\section{Proof of strict spreading in dimensions $1\le d\le4$}
\begin{proof}[Proof of Theorem~\ref{thm:monotone}]
With $A(s)$ and $D(s)$ as above, the derivative of $F$ is
\begin{equation}\label{eq:derivative}
 F'(s)=\frac{A(s)-s+4\alpha sA(s)}{(s+A(s))^3}.
\end{equation}
The combination that controls dilation is
\begin{align}
 4F(s)+dsF'(s)
 &=F(s)\left[
 4-d+\frac{2dA(s)(1+2\alpha s)}{s+A(s)}
 \right] \notag\\
 &=F(s)\left[
 4-d+\frac{2d\varepsilon^2e^{-2\alpha s}(1+2\alpha s)}
 {s+\varepsilon^2e^{-2\alpha s}}
 \right].\label{eq:virial_weight}
\end{align}
For $1\le d\le4$, the bracket is strictly positive at every
finite $s\ge0$. Accordingly,
\begin{equation}\label{eq:virial_sign}
 4F(s)+dsF'(s)\ge0\quad(s\ge0),
 \qquad 4F(s)+dsF'(s)>0\quad(s>0).
\end{equation}

To justify differentiating the nonlinear integral in
\eqref{eq:scale}, note that $sF'(s)$ is bounded. Indeed,
\eqref{eq:derivative} makes $sF'(s)$ continuous on
$[0,\infty)$ with limit zero at both $s=0$ and $s=\infty$.
Thus $4F+dsF'$ is bounded on $[0,\infty)$.
For $\lambda$ in any compact subinterval of $(0,\infty)$,
the derivative with respect to $\log\lambda$ of the
integrand $\lambda^4F(\lambda^d\rho)B$ is bounded in
absolute value by a constant times $B\in L^1(\R^d)$.
Dominated convergence therefore yields
\begin{align}
 \frac{\dd}{\dd\log\lambda}E(\psi_\lambda)
 &=2c_1\lambda^2\|\nabla\psi\|_{L^2}^2\notag\\
 &\quad+c_2\lambda^4\int_{\R^d}
 \bigl(4F(s)+dsF'(s)\bigr)B(y)\dd y,
 \qquad s=\lambda^d\rho(y).
 \label{eq:monotone_derivative}
\end{align}
By \eqref{eq:virial_sign}, the integral is nonnegative.
By \eqref{eq:positive}, $\|\nabla\psi\|_{L^2}^2>0$ for
nonzero $\psi$. Therefore
\begin{equation}\label{eq:positive_derivative}
 \boxed{\frac{\dd}{\dd\log\lambda}E(\psi_\lambda)
 \ge2c_1\lambda^2\|\nabla\psi\|_{L^2}^2>0.}
\end{equation}
Strict monotonicity and \eqref{eq:strict} follow.

For completeness, the dilation group is strongly continuous on
$H^2(\R^d)$:
\begin{equation}\label{eq:strong_cont}
 \lim_{\lambda\to1}\|\psi_\lambda-\psi\|_{H^2}=0.
\end{equation}
To see this, first use dominated convergence for functions in
$C_c^\infty(\R^d)$, and then pass to all of $H^2$ by density;
the dilation operators are uniformly bounded on $H^2$ for
$\lambda$ in a neighborhood of $1$, because their
$L^2$, first-derivative, and second-derivative norms scale
respectively as $1$, $\lambda$, and $\lambda^2$.

Now suppose that $\psi\in\A_m$ were a local constrained
$H^2$ minimizer. By \eqref{eq:strong_cont}, for $\lambda<1$
sufficiently close to $1$, $\psi_\lambda$ would lie in its
minimizing neighborhood. Yet \eqref{eq:mass_preserved} keeps
$\psi_\lambda\in\A_m$ while \eqref{eq:strict} implies
$E(\psi_\lambda)<E(\psi)$, a contradiction.
\end{proof}

\section{Interpretive limits}
\begin{remark}[Dimensions greater than four]
The proof of nonattainment uses only boundedness of $F$,
positivity of the Dirichlet term, and mass-preserving
spreading. It is therefore valid for every finite $d\ge1$.
For $d>4$, the factor $4-d$ in \eqref{eq:virial_weight}
can be negative. Thus the argument above does not establish
strict dilation monotonicity or the absence of \emph{local}
constrained minimizers in those dimensions.
\end{remark}

\begin{remark}[Stationary points and other models]
The absence of positive-mass minimizing states is an energetic
statement, not a nonexistence theorem for all stationary
saddles or a dynamical scattering theorem. Likewise, excluding
general constrained critical points via dilation requires
justifying the dilation as a differentiable tangent variation
in the selected function space. The result does not cover
additional focusing terms, external confinement, prescribed
boundaries, or other energy functionals with different signs.
\end{remark}

\section*{Source and reproducibility}
A related foundational audit and scope discussion is available at
\url{https://github.com/rickyjreyes/geometry_of_resonance/blob/main/docs/foundational/LOCALIZED_STATE_EXISTENCE.md}.

\end{document}
